\section{Introduction}
\label{sec:intro}

Frameworks for emergence invite a standard objection: that they merely redescribe. A theory that ``explains'' why a logical qubit is a robust object by renaming syndrome extraction as \emph{packaging} and decoding as \emph{lifting} has added vocabulary, not knowledge. Six Birds Theory (SBT)~\cite{sbt-f1,sbt-f3} is a framework in this genre, and its own foundational papers state the standard it must meet: the framework earns its keep only where it changes the \emph{audit layer}---only where its primitives compute something checkable about the domain, and where those computations can fail. This paper is an adversarial test of that standard, in the falsificationist sense~\cite{popper1959logic,mayo2018severe}, carried out on quantum error correction.

The test design is simple to state. We took SBT's six audit primitives---the existence/closure certificate, the adequacy residual, the predictive quotient (memory witness), the protocol trap, the shadow-price/budget audit, and transport---all defined on top of its packaging operator, and asked what each one \emph{must} become when the domain is a small stabilizer code under stochastic Pauli noise. We then implemented the resulting objects for five of the six---transport was scoped out before the build began and never implemented---in a deterministic simulation prototype whose theorem-grade quantities are computed in exact rational arithmetic, froze quantitative predictions about their behavior \emph{before} each scored run (for E1--E7 before implementation; for the two circuit-level experiments E8--E9 from a pilot on disjoint seeds, Section~\ref{sec:methods-registration}), and ran the experiments against those frozen bars. A failed prediction is a first-class result: it is reported with its frozen bar, its measured value, and a diagnostic reading, and it is never re-run with tuned parameters.

Stating the outcome honestly requires a split, and we make it here rather than in a limitations section. Four of the six primitives produced positively evidenced, machine-checkable audit objects: the adequacy residual (a priced answer to ``do the scheduled checks cover the logical degrees of freedom?''), the existence certificate (a typed status for whether a stable logical layer exists at all), the predictive-quotient memory witness (an exact certificate that histories indistinguishable by current observables differ in their declared futures, with the operational value of that memory measured separately), and the protocol-trap control (which ran correctly and falsified our own ``this is a schedule artifact'' hypothesis). A fifth primitive, the pricing audit, was implemented exactly and its machinery works, but every one of its registered predictions failed at the frozen tolerances---we report it as a negative. A sixth, transport/functoriality, was explicitly deferred before the build began and was never implemented. % C-01; counts F-147, F-148; split per REPORT.md §6
Across the original build, 25 predictions were frozen and 14 of them failed as frozen. % F-147, F-148 (C-02)

Why publish a prototype with a 56\% headline failure rate, rather than leave the failures in the file drawer~\cite{rosenthal1979filedrawer}? Because the failures are where the framework's value showed up. % F-147, F-148 (14/25)
The paper's centerpiece (Section~\ref{sec:arc}) is what happened to one of them. A deployable, syndrome-only drift-detection witness registered negative: it never detected the frozen drift injection within the tested grid (up to $N = 16000$ shots). Rather than tune it, we diagnosed it \emph{inside the framework}. Exact computation---including a Walsh--Hadamard inversion that recovers the full syndrome distribution in closed form---located four findings: two exact defects in the statistic, one mis-calibrated null confirmed by measurement, and one structural fact about the frozen test scenario itself (it sits, by construction, nearly on a minimum-weight logical-failure path, the scenario among the three tested where a syndrome-only method is at the greatest disadvantage). The corrected witness was then re-registered with new frozen predictions, including one \emph{predicting that the corrected witness would still not beat the baseline on the original scenario}. That predicted failure held, alongside a dramatic predicted success on an off-support scenario. % C-03
The correction then generalized to circuit-level noise (clearing its registered bar at the thinnest margin the test grid resolves---we say so every time) and to a closed decoder-recalibration loop on the one channel known by construction to drift, where witness-triggered recalibration landed about $0.0014$ above an oracle ceiling in post-drift logical error rate and beat a blind schedule matched to the \emph{same per-seed recalibration budget} by $0.0185$; whether it also beats a five-times-larger blind budget is unresolved at ten seeds. % F-091, F-092, F-094, F-095 (C-34, C-35, C-39)

To be equally plain about what this paper does \emph{not} claim: none of the individual QEC mechanisms is new. Detector-correlation calibration is standard experimental practice~\cite{fowler2014scalable,chen2021exponential,google2023suppressing}, decoder reweighting under drift exists~\cite{spitz2018adaptive,huo2017realtime,dgr2023reweighting,bhardwaj2025drifting}, and sequential change detection is seventy years old~\cite{page1954continuous}. No hardware was involved, no asymptotic threshold statements are made, and the residual certificate is explicitly relative to a declared estimator class (Section~\ref{sec:instantiation}). The contribution is the test itself and its outcome: when SBT's primitives were forced into contact with a real technical domain under pre-registration discipline and adversarial internal and external review, they generated checkable structure, genuine falsifications, and a repair loop that operated within the framework rather than around it.

The paper is organized so that a reader who has never heard of SBT can follow it. Section~\ref{sec:sbt} introduces the framework from scratch, in the concrete finite setting actually used. Section~\ref{sec:instantiation} maps each primitive to its QEC object and states the scope fences. Section~\ref{sec:methods} describes the evidence discipline: exactness, registered predictions, controls, and reproducibility. Sections~\ref{sec:results-certificates}--\ref{sec:arc} present results: the certificates on ground truth, the negative results as a taxonomy, and the diagnosis-and-repair arc. Sections~\ref{sec:related}--\ref{sec:discussion} place the work and state its limits, and Section~\ref{sec:repro} gives the exact environment and commands that regenerate every number byte-identically. Appendices carry the notation table, the complete prediction ledger, the controls ledger, the deviations log, and a formula summary.

% ===========================================================================
